%% P09 --- Eigenvalues, quadratic forms and positive definiteness
%%
%% Stub. The brief below is the contract for this program: what it must argue,
%% and what it must buy the reader. Writing the program means deleting the
%% \programstub{} block and replacing it with the program. If the brief turns
%% out to be wrong once the mathematics has been worked, change it and record
%% the contradiction in CLAUDE.md under *Resolved questions*.

\program{Eigenvalues, quadratic forms and positive definiteness}
\label{prog:P09}

\programstub{%
Argument: an eigenvector is a direction the matrix only stretches; a
symmetric matrix has a full orthogonal set of them (the spectral theorem,
stated and used, not proved); a quadratic form is a bowl, a saddle or a
ridge, and its eigenvalues say which. Payoff: a covariance matrix is
symmetric and positive semi-definite, so PCA is an eigen-decomposition and
not a black box; the spectral norm as the largest stretch, which is the
Lipschitz constant that bounds how much a layer can amplify; the Hessian's
eigenvalues are the shape of the loss basin, which is where \enquote{ravine}, \enquote{sharp minimum} and \enquote{the learning rate is too high} all become one statement ---
collected here and spent in P16 and P19.

\medskip
\textbf{Planned length:} 65 frames.
}
